\[ %%%%%%%%%%%%%%%%%%%%%%%%%%%%%%%%%%%%%%%%%%%%%%%%%%%%%%%%%%%%%%%%%%%%%%%%%%%%%%%% \subsection{Differentiable Physics and Hybrid Simulation} \label{sec:differentiable_physics} %%%%%%%%%%%%%%%%%%%%%%%%%%%%%%%%%%%%%%%%%%%%%%%%%%%%%%%%%%%%%%%%%%%%%%%%%%%%%%%% Differentiable physics has emerged as an important research direction aimed at combining the predictive capabilities of machine learning with the accuracy and interpretability of physics-based simulation. The central idea of differentiable simulation is to construct computational models in which physical simulators become differentiable components of a larger optimization and learning pipeline. This enables gradients to propagate through numerical simulation procedures, allowing machine learning algorithms to directly interact with physical models. In conventional scientific computing, numerical solvers are typically treated as independent computational tools that generate solutions for predefined physical models. Although these methods provide high accuracy and strong physical guarantees, they often require substantial computational resources and may lack adaptability when applied to inverse problems, parameter identification, optimization, and real-time control. Differentiable physics introduces a new paradigm in which simulation and learning become tightly integrated. A general differentiable simulation framework can be represented as: \begin{equation} x_{t+1}=S_\phi(x_t,u_t), \end{equation} where $S_\phi$ denotes a differentiable simulator parameterized by physical or computational parameters $\phi$, $x_t$ represents the system state, and $u_t$ denotes external inputs or control variables. By computing gradients through $S_\phi$, optimization algorithms can adjust parameters, infer unknown physical quantities, or train neural components embedded within the simulation pipeline. This approach has enabled significant advances in several scientific domains, including computational fluid dynamics, robotics, molecular simulation, deformable object modeling, and physical system identification. Differentiable simulators allow neural networks to complement traditional numerical methods by learning unresolved dynamics, correcting model discrepancies, or accelerating expensive computational components. Hybrid simulation approaches extend this concept by combining analytical physical models with data-driven learning modules. Instead of replacing physics-based solvers entirely, hybrid models assign different computational responsibilities to physical and neural components. For example, known physical laws may describe large-scale system behavior, while neural networks approximate unknown constitutive relations, unresolved sub-grid phenomena, or computationally expensive operators. Neural ordinary differential equations (Neural ODEs) represent another important development within this direction. By parameterizing continuous-time dynamics using neural networks, \begin{equation} \frac{dx}{dt}=f_\theta(x,t), \end{equation} Neural ODEs provide a flexible framework for learning unknown dynamical systems while preserving continuous temporal evolution. Extensions to neural differential equations, stochastic differential equations, and Hamiltonian neural networks have further expanded the ability of machine learning models to represent complex physical processes. Despite these achievements, differentiable physics and hybrid simulation approaches face several challenges. First, many differentiable simulators depend on accurate underlying physical models. When governing equations are incomplete, unknown, or computationally prohibitive, the effectiveness of purely differentiable simulation may be limited. Second, gradient computation through complex numerical procedures can introduce significant computational overhead. Third, maintaining long-horizon stability remains challenging when learned components are used recursively within dynamical simulations. Furthermore, many hybrid approaches rely on predefined simulation structures and do not naturally generalize to systems represented by changing topologies, irregular interactions, or heterogeneous components. Real-world physical systems often require flexible representations capable of adapting to evolving spatial relationships and multi-scale interactions. The PHYSGEN framework builds upon the principles of differentiable physics while extending them toward graph-based physics-guided learning. Instead of embedding neural components into a fixed numerical simulator, PHYSGEN treats physical systems as dynamic interaction graphs and introduces differentiable physics guidance directly into graph evolution, message passing, and latent state propagation. This perspective enables a flexible combination of data-driven learning and physical reasoning. Analytical models, numerical simulations, and observational data can all contribute to the learning process while the graph architecture provides a universal representation for complex interacting systems. Therefore, PHYSGEN can be considered a physics-guided hybrid computational framework that combines the adaptability of machine learning with the structural consistency of physical simulation, providing a foundation for stable long-horizon prediction and future digital twin applications. \] 